\subsection{Señales sintéticas de ECG}

Similar a los ejercicios anteriores, empezaremos importando las bibliotecas 
necesarias para trabajar con señales unidimensionales y aplicar los filtros 
correspondientes.

\begin{lstlisting}[language = Python, caption = {Carga de bibliotecas para procesamiento de señales.}]
# Bibliotecas

import numpy as np
import matplotlib.pyplot as plt
from scipy.signal import butter, filtfilt
\end{lstlisting}

\subsubsection{Generación de la señal original}

Para crear nuestra señal sintética, utilizaremos la suma de funciones gaussianas. 
Cada función simulará una parte del complejo PQRST del latido cardíaco ajustando 
su amplitud, centro y ancho. Posteriormente, crearemos una señal continua sumando 
varios de estos latidos a lo largo de un vector de tiempo.

\begin{lstlisting}[language = Python, caption = {Generación de la señal sintética de ECG.}]
frecuencia = 500  # Frecuencia de muestreo 
t = np.arange(0, 3, 1/frecuencia)  # Vector de tiempo de 3 segundos

# Utilizamos funciones Gaussianas para simular el PQRST
def simular_latido(t_local):
    ondas = [
        (0.15, -0.2, 0.02),   # Onda P
        (-0.15, -0.05, 0.01), # Onda Q
        (1.0, 0.0, 0.015),    # Onda R
        (-0.25, 0.05, 0.015), # Onda S
        (0.3, 0.3, 0.04)      # Onda T
    ]
    latido = np.zeros_like(t_local)
    for a, mu, sigma in ondas:
        latido += a * np.exp(-0.5 * ((t_local - mu) / sigma)**2)
    return latido

# Creamos una señal vacía y le sumamos latidos cada 1 segundo
ecg_original = np.zeros_like(t)
for i in range(4): # Latidos en los segundos 0, 1, 2 y 3
    ecg_original += simular_latido(t - i)
\end{lstlisting}

\begin{figure}[H]
    \centering
    \includegraphics[width=0.8\textwidth]{ecg_original.png}
    \caption{Señal sintética de ECG original.}
    \label{fig:ecg_original}
\end{figure}

\subsubsection{Adición de ruido a la señal}

Para emular un escenario realista y poner a prueba los filtros, agregaremos 
dos tipos de ruido a nuestra señal original: una interferencia eléctrica 
simulada con una onda senoidal de 60 Hz, y ruido aleatorio (ruido gaussiano) 
para emular la estática de los sensores.

\begin{lstlisting}[language=Python, caption={Adición de ruido de 60 Hz y ruido gaussiano.}]
ruido_60hz = 0.2 * np.sin(2 * np.pi * 60 * t)       # Interferencia electrica (60 Hz)
ruido_gauss = np.random.normal(0, 0.08, len(t))     # Ruido aleatorio (Gaussiano)

ecg_ruidoso = ecg_original + ruido_60hz + ruido_gauss
\end{lstlisting}

\begin{figure}[H]
    \centering
    \includegraphics[width=0.8\textwidth]{ecg_ruido.png}
    \caption{Señal de ECG con ruido.}
    \label{fig:ecg_ruido}
\end{figure}

\subsubsection{Transformada de Fourier y filtrado en frecuencia}

Para analizar las componentes de frecuencia de la señal contaminada, utilizaremos 
la Transformada Rápida de Fourier (FFT). Posteriormente, diseñaremos un filtro ideal 
en el dominio de la frecuencia estableciendo a cero todas las frecuencias superiores 
a 40 Hz, y regresaremos al dominio del tiempo con la transformada inversa (IFFT).

\begin{lstlisting}[language=Python, caption={Aplicación de la FFT y filtrado en el dominio de la frecuencia.}]
N = len(t)
freqs = np.fft.fftfreq(N, 1/frecuencia)
fft_ruidoso = np.fft.fft(ecg_ruidoso)

# Hacemos una copia y establecemos a cero las frecuencias mayores a 40 Hz
fft_filtrado = fft_ruidoso.copy()
fft_filtrado[np.abs(freqs) > 40] = 0  

# Transformada inversa (IFFT) para volver al dominio del tiempo
ecg_filtrado_freq = np.real(np.fft.ifft(fft_filtrado))
\end{lstlisting}

\begin{figure}[H]
    \centering
    \includegraphics[width=0.8\textwidth]{ecg_espectro.png}
    \caption{Espectro de frecuencia de la señal con ruido.}
    \label{fig:ecg_espectro}
\end{figure}

\subsubsection{Filtrado en el dominio del tiempo}

De manera opcional y para tener un punto de comparación, se diseñó un filtro 
pasabajas de Butterworth de orden 4 con una frecuencia de corte de 40 Hz.
 Para evitar el desfasamiento temporal de la señal, se aplicó el filtro 
 utilizando la función \texttt{filtfilt}.

\begin{lstlisting}[language=Python, caption={Filtro Butterworth pasabajas en el dominio del tiempo.}]
# Filtro Butterworth pasa bajas de orden 4, frecuencia de corte 40 Hz
b, a = butter(4, 40, fs=frecuencia, btype='low')

# Usamos filtfilt para no generar desfasaje temporal
ecg_filtrado_tiempo = filtfilt(b, a, ecg_ruidoso)
\end{lstlisting}

\subsubsection{Visualización y comparación de resultados}

Para finalizar esta sección, se preparará la visualización de los resultados 
superponiendo la señal original, la señal con ruido y las dos señales filtradas 
para observar la efectividad de los métodos aplicados.

\begin{lstlisting}[language=Python, caption={Comparación y superposición de las señales de ECG.}]
plt.figure(figsize = (12, 6))

plt.plot(t, ecg_ruidoso, label = 'Señal con Ruido', color = 'gray', alpha = 0.3)
plt.plot(t, ecg_original, label = 'Original', color = 'green', linewidth = 1.5, alpha = 0.8)
plt.plot(t, ecg_filtrado_freq, label = 'Filtrada (Frecuencia)', color = 'red', linestyle = 'dashed', linewidth = 2, alpha = 0.6)
plt.plot(t, ecg_filtrado_tiempo, label = 'Filtrada (Tiempo)', color = 'blue', linestyle = 'dotted', linewidth = 2)

plt.title('Comparación y Superposición de Señales')
plt.xlabel('Tiempo (s)')
plt.ylabel('Amplitud')
plt.xlim(0, 3)
plt.legend(loc = 'upper right')
plt.grid(True, alpha = 0.3, linestyle = '--')
plt.show()
\end{lstlisting}

\begin{figure}[H]
    \centering
    \includegraphics[width=\textwidth]{ecg_comparacion.png}
    \caption{Comparación de la efectividad de los filtros en la señal de ECG.}
    \label{fig:ecg_comparacion}
\end{figure}